% Generated draft: requires mathematical and source review.
\documentclass[11pt]{article}
\usepackage[margin=25mm]{geometry}
\usepackage{fontspec}
\setmainfont{Latin Modern Roman}
\usepackage{amsmath,amssymb,mathtools,unicode-math}
\setmathfont{Latin Modern Math}
\usepackage{enumitem,needspace,xurl,hyperref}
\setlength{\parindent}{0pt}
\setlength{\parskip}{5pt}
\setlength{\emergencystretch}{3em}
\newcommand{\N}{\mathbb{N}}
\newcommand{\Z}{\mathbb{Z}}
\newcommand{\Q}{\mathbb{Q}}
\newcommand{\R}{\mathbb{R}}
\newcommand{\C}{\mathbb{C}}
\newcommand{\F}{\mathbb{F}}
\newcommand{\A}{\mathbb{A}}
\newcommand{\PP}{\mathbb{P}}
\newcommand{\E}{\mathbb{E}}
\newcommand{\eps}{\varepsilon}
\newcommand{\dd}{\,\mathrm{d}}
\begin{document}
\section*{UnsolvedMath ID: 10600028}
\textbf{Source problem code:} catalogue entry
\subsection{Definitions and mathematical statement}
Fix a convention for biquandles, namely the algebraic crossing operations and invertibility and Reidemeister axioms used by the source. The supplied excerpt does not list these axioms. In that category, a free biquandle on a set $X$ is a biquandle $F(X)$ with a map $\eta:X\to F(X)$ such that, for every biquandle $B$ and every map $f:X\to B$, there exists a unique biquandle homomorphism $\overline f:F(X)\to B$ with $\overline f\eta=f$.

Problem-list item 28 retains three questions:
\begin{enumerate}
\item Give a descriptive representation of the free biquandle, in addition to its abstract universal property or formal generators-and-relations description.
\item Give a topological explanation of the fundamental biquandle of a knot. In the usual diagrammatic construction, generators correspond to semiarcs and relations encode the crossing operations; the request is for a topological understanding of this object, with conventions compatible with that presentation.
\item Determine whether there is a useful notion of a free partial biquandle. The source expresses the tentative expectation ``probably not,'' rather than a proved nonexistence assertion.
\end{enumerate}
The meanings of partial biquandle and descriptive representation, the generator scope, and the desired intrinsic topological construction are not specified. Thus the last question has no fixed universal property until a category of partial objects and its morphisms is defined. The excerpt also leaves the intended classical or virtual knot scope and orientation conventions unstated.
\subsection{Short English statement}
Find a concrete description of free biquandles and a topological interpretation of knot biquandles, and investigate whether an appropriate free partial biquandle can exist. The source's negative expectation for the last question is tentative.
\subsection{Sources}
\begin{itemize}
\item Local registry: \url{data/unsolvedmath/problems.json}, permanent UnsolvedMath ID $10600028$, exact source code \texttt{catalogue entry}, ``Virtual-knot problem 28 --- Biquandles.''
\item \emph{Unsolved Problems in Virtual Knot Theory and Combinatorial Knot Theory}, revised survey, problem-list item 28: \url{https://arxiv.org/abs/1409.2823}. Supplied access date: 2026-07-29; extraction: source-tex. All three parts of this compound item are retained.
\item The registry mentions general biquandle constructions with the incomplete identifier \texttt{arXiv:1908.xxxx} (2019). This is not a usable bibliographic citation; no precise additional reference or individual attribution is supplied.
\end{itemize}
\end{document}
